\documentclass[12pt]{article}
\usepackage[margin=2.5cm]{geometry}
\usepackage{parskip}
\usepackage{times}
\usepackage{xcolor}
\usepackage{hyperref}

\begin{document}

\begin{flushright}
\textbf{Yan Xiao} (Corresponding Author)\\
Research Center of Satellite Technology\\
State Key Laboratory of Micro-Spacecraft Rapid Design\\
and Intelligent Cluster\\
Harbin Institute of Technology\\
Harbin 150001, China\\
\href{mailto:xiaoy@hit.edu.cn}{xiaoy@hit.edu.cn}\\[6pt]
June 10, 2026
\end{flushright}

\bigskip

Dr.\ Maruthi Akella\\
Associate Editor\\
\textit{Journal of Guidance, Control, and Dynamics}\\
American Institute of Aeronautics and Astronautics

\bigskip

\textbf{Re: Manuscript \#2026-04-G010165 (Revision)}\\[4pt]
\textbf{``Exact Prescribed-Time Prescribed-Accuracy Control for Spacecraft
under Dual Reaction Wheel Saturation''}

\bigskip

Dear Dr.\ Akella,

We sincerely thank you for the editorial screening. Following your
recommendation, we have revised the manuscript from a Full Paper to a
Technical Note. The abstract has been removed, and the body text has been
substantially condensed to approximately 2500 words while preserving the core
contributions.

The revised manuscript focuses on two specific contributions that distinguish
it from existing methods.

\textbf{Exact prescribed-time convergence.}
Existing PTC methods are inherently conservative, converging well before the
prescribed terminal time $T_f$. The proposed two-loop architecture eliminates
this conservatism: the outer loop generates a reference trajectory that reaches
the target exactly at $T_f$, while the inner loop drives the tracking error to
within a prescribed accuracy despite uncertainties and disturbances.

\textbf{Unified treatment of dual saturation.}
Previous methods address torque saturation but largely neglect wheel
angular-momentum saturation. By enforcing both constraints,
the proposed framework guarantees that dual saturation is strictly satisfied
throughout the maneuver. Moreover, it provides an analytical estimate of the
minimum feasible maneuver time $T_{f,\min}$, ensuring that the prescribed
terminal time is physically realizable.

We hope the revised manuscript is now suitable for consideration as a
Technical Note in JGCD. Should you have any questions, please do not hesitate
to contact the corresponding author.

Yours sincerely,

\bigskip\bigskip

Zaihua Liu, Yan Xiao, Dong Ye, Haiyue Yang, and Zhaowei Sun

\end{document}
